\documentclass[aspectratio=169,11pt]{beamer}
\usetheme{default}
\setbeamertemplate{navigation symbols}{}
\setbeamertemplate{footline}{%
  \hfill\insertframenumber/\inserttotalframenumber\hspace*{4pt}\vspace{2pt}%
}
\usepackage{lmodern}
\usepackage[most]{tcolorbox}
\usepackage{listings}
\usepackage{tikz}
\usetikzlibrary{arrows.meta,positioning,shapes.geometric,fit,backgrounds}
\usepackage{booktabs}
\usepackage{hyperref}

\definecolor{lcgreen}{HTML}{2E7D32}
\definecolor{lcblue}{HTML}{1565C0}
\definecolor{lcorange}{HTML}{EF6C00}
\definecolor{lcred}{HTML}{C62828}
\definecolor{codebg}{HTML}{F5F5F5}
\definecolor{docgray}{gray}{0.55}

\newtcolorbox{conceptbox}[1][]{colback=yellow!20,colframe=yellow!60!black,title=#1,fonttitle=\bfseries,top=1pt,bottom=1pt,left=3pt,right=3pt}
\newtcolorbox{warningbox}[1][]{colback=red!15,colframe=red!60!black,title=#1,fonttitle=\bfseries,top=1pt,bottom=1pt,left=3pt,right=3pt}
\newtcolorbox{tipbox}[1][]{colback=green!10,colframe=green!50!black,title=#1,fonttitle=\bfseries,top=1pt,bottom=1pt,left=3pt,right=3pt}
\newtcolorbox{codebox}[1][]{colback=blue!10,colframe=blue!40!black,title=#1,fonttitle=\bfseries,top=1pt,bottom=1pt,left=3pt,right=3pt}

\lstset{
  basicstyle=\ttfamily\tiny,
  backgroundcolor=\color{codebg},
  breaklines=true,
  frame=single,
  rulecolor=\color{blue!40!black},
  keywordstyle=\color{lcblue}\bfseries,
  stringstyle=\color{lcgreen},
  commentstyle=\color{docgray}\itshape,
  language=Python,
  showstringspaces=false,
  tabsize=2,
  xleftmargin=2pt,
  xrightmargin=2pt,
  aboveskip=2pt,
  belowskip=2pt,
}

\newcommand{\docref}[1]{{\scriptsize\color{docgray}[Docs: #1]}}

\title{Section 11: LangGraph Fundamentals}
\subtitle{Graph-Based Orchestration for Stateful Agents}
\author{LangChain v1.2.x}
\date{March 2026}

\begin{document}

% ============================================================================
% SLIDE 1: Title Slide
% ============================================================================
\frame{\titlepage}

% ============================================================================
% SLIDE 2: Section Overview
% ============================================================================
\begin{frame}[fragile]
\frametitle{What We'll Cover}
\begin{itemize}
    \setlength{\itemsep}{1pt}
  \item Mental model: LangGraph as a graph-based orchestration framework
  \item LangGraph vs LangChain: relationships and use cases
  \item Core concepts: StateGraph, nodes, edges, START/END
  \item Building stateful, looping agents with cycles
  \item Checkpointing and persistence for resilience
  \item Human-in-the-loop interrupts and resumption
  \item Streaming, subgraphs, and dynamic graphs
  \item Patterns, debugging, common pitfalls
\end{itemize}
\vspace{10pt}
\textbf{Goal:} Master graph-based orchestration for complex agent workflows.
\end{frame}

% ============================================================================
% SLIDE 3: Mental Model
% ============================================================================
\begin{frame}
\frametitle{Mental Model: LangGraph as Graph Orchestration}
\begin{itemize}
    \setlength{\itemsep}{1pt}
  \item \textbf{LangGraph:} A graph-based framework for orchestrating multi-step agent workflows
  \item Nodes = functions that read/write \textit{state}
  \item Edges = transitions between nodes (conditional or fixed)
  \item \textbf{State:} Shared data structure (TypedDict) passed through the graph
  \item Graph \textit{runs} in discrete steps: execute a node, update state, route to next node
  \item \textbf{Stateful:} Unlike chains, graphs maintain and evolve state throughout execution
  \item \textbf{Looping:} Supports cycles, enabling iterative reasoning and refinement
\end{itemize}
\vspace{10pt}
\begin{center}
  \textit{Agents as Graphs: think in terms of state transitions, not function composition}
\end{center}
\end{frame}

% ============================================================================
% SLIDE 4: Why LangGraph? Problems It Solves
% ============================================================================
\begin{frame}
\frametitle{Why LangGraph? Problems It Solves}
\begin{itemize}
    \setlength{\itemsep}{1pt}
  \item \textbf{Cycles \& Loops:} Chains are DAGs (acyclic). Agents need feedback loops for reflection, retries, iteration.
  \item \textbf{Explicit State:} Chains pass one output as next input. LangGraph manages rich, multi-field state.
  \item \textbf{Conditional Routing:} Complex branching logic needs explicit graph structure, not nested if/else.
  \item \textbf{Persistence \& Recovery:} Checkpointing state at each step enables interrupts, resumption, and audit trails.
  \item \textbf{Human-in-the-Loop:} Pause execution at specific nodes, let humans review/modify, resume seamlessly.
  \item \textbf{Composability:} Subgraphs enable modular agent designs.
\end{itemize}
\end{frame}

% ============================================================================
% SLIDE 5: LangGraph vs LangChain
% ============================================================================
\begin{frame}
\frametitle{LangGraph vs LangChain}
\vspace{10pt}
\small
\begin{tabular}{l|c|c}
\toprule
\textbf{Feature} & \textbf{LangChain} & \textbf{LangGraph} \\
\midrule
Structure & DAG (Acyclic Chain) & General Graph (with cycles) \\
State Management & Linear pass-through & Explicit TypedDict \\
Loops & Not directly supported & First-class citizen \\
Conditional Routes & Limited & Full branching/conditions \\
Checkpointing & Basic & Rich, fine-grained \\
Human Interrupts & No & Yes (interrupt\_before, interrupt\_after) \\
Ideal For & Simple pipelines & Complex agents, multi-step workflows \\
\bottomrule
\end{tabular}

\vspace{15pt}
\textit{LangChain for chains; LangGraph for agents.}
\end{frame}

% ============================================================================
% SLIDE 6: Core Concept - StateGraph
% ============================================================================
\begin{frame}[fragile]
\frametitle{Core Concept: StateGraph}
\begin{itemize}
    \setlength{\itemsep}{1pt}
  \item \textbf{StateGraph:} A class that builds a graph with explicit state
  \item State is defined as a \texttt{TypedDict} (modern Python 3.8+)
  \item StateGraph constructor takes the state class
  \item You add nodes, edges, and configure entrypoint
\end{itemize}

\vspace{10pt}
\begin{lstlisting}
from typing_extensions import TypedDict
from langgraph.graph import StateGraph

class AgentState(TypedDict):
    input: str
    output: str
    messages: list

graph = StateGraph(AgentState)
# Add nodes, edges, compile...
\end{lstlisting}

\vspace{10pt}
\textbf{Key:} State is the communication medium between nodes.
\end{frame}

% ============================================================================
% SLIDE 7: Nodes - Reading and Writing State
% ============================================================================
\begin{frame}[fragile]
\frametitle{Nodes: Functions That Read and Write State}
\begin{itemize}
    \setlength{\itemsep}{1pt}
  \item \textbf{Node:} A Python function that receives state, performs work, returns updated state
  \item Signature: \texttt{def node\_name(state: StateType) -> dict}
  \item \textbf{Return value:} A dict with keys to update in state (merged, not replaced)
  \item Nodes can call LLMs, tools, compute, or any Python code
\end{itemize}

\vspace{10pt}
\begin{lstlisting}
def process_node(state: AgentState) -> dict:
    """A node that processes the input."""
    result = state["input"].upper()  # Read state
    return {"output": result}  # Write state (merged)

# Add node to graph
graph.add_node("process", process_node)
\end{lstlisting}

\vspace{10pt}
\textit{Nodes are the "verbs" of your agent graph.}
\end{frame}

% ============================================================================
% SLIDE 8: Edges - Connecting Nodes
% ============================================================================
\begin{frame}[fragile]
\frametitle{Edges: Connecting Nodes}
\begin{itemize}
    \setlength{\itemsep}{1pt}
  \item \textbf{Edge:} A transition from one node to another
  \item \textbf{Normal edge:} \texttt{add\_edge(source, target)} — always routes from source to target
  \item \textbf{Conditional edge:} \texttt{add\_conditional\_edges(source, routing\_func, mapping)}
  \item Routing function receives state, returns next node name
\end{itemize}

\vspace{10pt}
\begin{lstlisting}
# Normal edge: always go from "node_a" to "node_b"
graph.add_edge("node_a", "node_b")

# Conditional edge: routing function decides next node
def route(state: AgentState) -> str:
    if state["output"] == "DONE":
        return "end"
    return "loop"

graph.add_conditional_edges("process", route,
    {"end": END, "loop": "process"})
\end{lstlisting}
\end{frame}

% ============================================================================
% SLIDE 9: START and END Special Nodes
% ============================================================================
\begin{frame}[fragile]
\frametitle{START and END: Special Nodes}
\begin{itemize}
    \setlength{\itemsep}{1pt}
  \item \textbf{START:} Virtual node representing graph entry point
  \item \textbf{END:} Virtual node representing graph exit
  \item Import: \texttt{from langgraph.graph import START, END}
  \item Always have at least one edge from START and one to END
\end{itemize}

\vspace{10pt}
\begin{lstlisting}
from langgraph.graph import START, END

# Route from START to first node
graph.add_edge(START, "input_node")

# Route to END when done
graph.add_edge("final_node", END)
\end{lstlisting}

\vspace{15pt}
\begin{tipbox}[Graph Structure]
Every graph must have a clear entry point (START) and exit (END).
\end{tipbox}
\end{frame}

% ============================================================================
% SLIDE 10: Building Your First Graph - Simple Chatbot
% ============================================================================
\begin{frame}[fragile]
\frametitle{Building Your First Graph: Simple Chatbot}
\begin{lstlisting}
from typing_extensions import TypedDict
from langgraph.graph import StateGraph, START, END

class ChatState(TypedDict):
    user_input: str
    response: str

def chat_node(state: ChatState) -> dict:
    """Simple echo chatbot."""
    response = f"You said: {state['user_input']}"
    return {"response": response}

graph = StateGraph(ChatState)
graph.add_node("chat", chat_node)
graph.add_edge(START, "chat")
graph.add_edge("chat", END)

compiled = graph.compile()
result = compiled.invoke({"user_input": "Hello"})
print(result["response"])  # "You said: Hello"
\end{lstlisting}
\end{frame}

% ============================================================================
% SLIDE 11: State Management - How State Flows
% ============================================================================
\begin{frame}
\frametitle{State Management: How State Flows Through the Graph}
\begin{itemize}
    \setlength{\itemsep}{1pt}
  \item State is initialized when you call \texttt{invoke(initial\_state)} or \texttt{stream(initial\_state)}
  \item Each node receives the current state (all fields)
  \item Node returns a dict with updates; these are \textbf{merged} into state
  \item Merging: all keys in state persist unless overwritten
  \item After node completes, routing logic executes (conditional edges)
  \item Graph continues until it reaches END
\end{itemize}

\vspace{15pt}
\begin{conceptbox}[State as Shared Memory]
State is not replaced between nodes—it's accumulated. Each node updates only the fields it needs to.
\end{conceptbox}
\end{frame}

% ============================================================================
% SLIDE 12: Conditional Routing with add_conditional_edges
% ============================================================================
\begin{frame}[fragile]
\frametitle{Conditional Routing: add\_conditional\_edges}
\begin{lstlisting}
def routing_decision(state: ChatState) -> str:
    """Decide next node based on state."""
    if "error" in state.get("response", "").lower():
        return "retry"
    return "end"

graph.add_conditional_edges(
    "process_node",
    routing_decision,
    {
        "retry": "process_node",    # Loop back
        "end": END                   # Exit
    }
)
\end{lstlisting}

\vspace{10pt}
\textbf{Key points:}
\begin{itemize}
    \setlength{\itemsep}{1pt}
  \item Routing function receives state, returns string (next node name)
  \item Mapping dict maps return value to actual node/END
  \item Enables loops, retries, branching based on runtime conditions
\end{itemize}
\end{frame}

% ============================================================================
% SLIDE 13: TikZ Diagram - Multi-Step Agent Graph
% ============================================================================
\begin{frame}
\frametitle{Multi-Step Agent Graph with Branching}
\vspace{-0.6cm}
\begin{center}
\begin{tikzpicture}[
  node distance=1cm,
  every node/.style={draw, rounded rectangle, minimum width=0.9cm, minimum height=0.35cm, font=\tiny},
  >=Stealth,
  scale=0.5
]

% Nodes
\node (start) [fill=lcgreen!70, text=white] {Start};
\node (input) [below=of start, fill=lcblue!70, text=white] {Input};
\node (think) [below=of input, fill=lcblue!70, text=white] {Think};
\node (decide) [below=of think, fill=lcorange!70, text=white] {Decide};
\node (act) [below left=0.8cm and 0.8cm of decide, fill=lcblue!70, text=white] {Act};
\node (refine) [below right=0.8cm and 0.8cm of decide, fill=lcblue!70, text=white] {Refine};
\node (end) [below right=1cm and 0.2cm of act, fill=lcred!70, text=white] {End};

% Edges
\draw [->, thick] (start) -- (input);
\draw [->, thick] (input) -- (think);
\draw [->, thick] (think) -- (decide);
\draw [->, thick] (decide) -- (act);
\draw [->, thick] (decide) -- (refine);
\draw [->, thick] (act) -- (end);
\draw [->, thick] (refine) -- (end);
\draw [<->, thick, dashed] (think) edge[bend left=40] (act);

\end{tikzpicture}
\end{center}
\vspace{-0.5cm}
\tiny
Feedback loop: Think $\to$ Act $\to$ Think
\end{frame}

% ============================================================================
% SLIDE 14: Cycles in Graphs - Enabling Loops
% ============================================================================
\begin{frame}
\frametitle{Cycles in Graphs: How LangGraph Enables Loops}
\begin{itemize}
    \setlength{\itemsep}{1pt}
  \item \textbf{LangChain chains:} DAGs (directed acyclic graphs) — no loops allowed
  \item \textbf{LangGraph:} Allows cycles — nodes can route back to earlier nodes
  \item \textbf{Loops enable:}
    \begin{itemize}
    \setlength{\itemsep}{1pt}
      \item Iterative reasoning: "think -> act -> observe -> think again"
      \item ReAct pattern: think, act, reflect, repeat
      \item Error recovery: try -> fail -> retry
      \item Refinement: generate -> evaluate -> improve
    \end{itemize}
  \item Cycles are \textit{explicit} — no hidden magic, you control when to loop
\end{itemize}

\vspace{15pt}
\begin{warningbox}[Infinite Loops]
Always ensure your routing logic has an exit condition, or set a max recursion depth.
\end{warningbox}
\end{frame}

% ============================================================================
% SLIDE 15: Checkpointing - MemorySaver
% ============================================================================
\begin{frame}[fragile]
\frametitle{Checkpointing: Saving and Restoring Graph State}
\begin{itemize}
    \setlength{\itemsep}{1pt}
  \item \textbf{Checkpointing:} Saving state after each node executes
  \item \textbf{MemorySaver:} In-memory checkpoint storage (for development)
\end{itemize}

\vspace{10pt}
\begin{lstlisting}
from langgraph.checkpoint.memory import MemorySaver

# Create checkpointer
checkpointer = MemorySaver()

# Compile with checkpointer
compiled = graph.compile(checkpointer=checkpointer)

# Invoke with thread_id to enable checkpointing
result = compiled.invoke(
    {"user_input": "Hello"},
    config={"configurable": {"thread_id": "abc123"}}
)
\end{lstlisting}

\vspace{10pt}
\textbf{Benefits:}
\begin{itemize}
    \setlength{\itemsep}{1pt}
  \item Resume execution from last checkpoint
  \item Inspect state at any step
  \item Implement human-in-the-loop workflows
\end{itemize}
\end{frame}

% ============================================================================
% SLIDE 16: Persistence Backends - SQLite and PostgreSQL
% ============================================================================
\begin{frame}[fragile]
\frametitle{Persistence Backends: SQLite and PostgreSQL}
\begin{lstlisting}
from langgraph.checkpoint.sqlite import SqliteSaver

# SQLite persistence
checkpointer = SqliteSaver(conn="langraph_db.sqlite")
compiled = graph.compile(checkpointer=checkpointer)
\end{lstlisting}

\vspace{10pt}
\begin{lstlisting}
from langgraph.checkpoint.postgres import PostgresSaver

# PostgreSQL persistence
checkpointer = PostgresSaver(
    conn_string="postgresql://user:pw@localhost:5432/db"
)
compiled = graph.compile(checkpointer=checkpointer)
\end{lstlisting}

\vspace{10pt}
\textbf{Key points:}
\begin{itemize}
    \setlength{\itemsep}{1pt}
  \item SQLite: Good for single-machine, development
  \item PostgreSQL: For distributed, production systems
  \item Config: Always use \texttt{config=\{"configurable": \{"thread\_id": id\}\}}
\end{itemize}
\end{frame}

% ============================================================================
% SLIDE 17: Human-in-the-Loop - Interrupts
% ============================================================================
\begin{frame}[fragile]
\frametitle{Human-in-the-Loop: interrupt\_before and interrupt\_after}
\begin{lstlisting}
# Pause BEFORE executing a node
compiled = graph.compile(
    checkpointer=checkpointer,
    interrupt_before=["approval_node"]
)

# Pause AFTER executing a node
compiled = graph.compile(
    checkpointer=checkpointer,
    interrupt_after=["llm_call_node"]
)
\end{lstlisting}

\vspace{10pt}
\textbf{Workflow:}
\begin{enumerate}
    \setlength{\itemsep}{1pt}
  \item Call \texttt{invoke()} or \texttt{stream()} — execution pauses at interrupt point
  \item Check state: \texttt{compiled.get\_state(config)}
  \item Inspect results, let human decide
  \item Optionally modify state
  \item Resume: \texttt{compiled.invoke(None, config)}
\end{enumerate}

\vspace{10pt}
Perfect for approval workflows, manual reviews, user feedback loops.
\end{frame}

% ============================================================================
% SLIDE 18: Resuming After Interrupts
% ============================================================================
\begin{frame}[fragile]
\frametitle{Resuming After Interrupts}
\begin{lstlisting}
# First execution - pauses at "approval"
config = {"configurable": {"thread_id": "user_session"}}
compiled.invoke({"task": "process"}, config)

# Inspect state at interrupt point
state = compiled.get_state(config)
print(f"Pending input: {state.values}")

# Modify state if needed
updated_state = {
    "task": "process",
    "approved": True,
    "notes": "Looks good"
}

# Resume execution
result = compiled.invoke(updated_state, config)
\end{lstlisting}

\vspace{10pt}
\textit{State persists; execution picks up where it left off.}
\end{frame}

% ============================================================================
% SLIDE 19: Streaming Graph Execution
% ============================================================================
\begin{frame}[fragile]
\frametitle{Streaming Graph Execution: stream() and astream\_events()}
\begin{itemize}
    \setlength{\itemsep}{1pt}
  \item \texttt{invoke()}: Returns final state only
  \item \texttt{stream()}: Yields state after each node (useful for long operations)
  \item \texttt{astream\_events()}: Async stream of intermediate events (for real-time updates)
\end{itemize}

\vspace{10pt}
\begin{lstlisting}
# Stream states after each node
for snapshot in compiled.stream({"input": "..."}, config):
    node_name, node_state = snapshot
    print(f"Node '{node_name}' completed")
    print(f"State: {node_state}")

# Async streaming
async for event in compiled.astream_events(input, config):
    print(f"Event type: {event['event']}, data: {event['data']}")
\end{lstlisting}
\end{frame}

% ============================================================================
% SLIDE 20: Subgraphs - Composing Graphs
% ============================================================================
\begin{frame}[fragile]
\frametitle{Subgraphs: Composing Graphs Within Graphs}
\begin{itemize}
    \setlength{\itemsep}{1pt}
  \item \textbf{Subgraph:} A compiled graph used as a node in another graph
  \item Enables modular, reusable agent components
  \item Subgraph state must align with parent state (or use adapters)
\end{itemize}

\vspace{10pt}
\begin{lstlisting}
# Define a sub-graph
sub_graph = StateGraph(AgentState).compile()

# Use it as a node in parent graph
parent_graph = StateGraph(AgentState)
parent_graph.add_node("subprocess", sub_graph)
parent_graph.add_edge(START, "subprocess")
parent_graph.add_edge("subprocess", END)

compiled = parent_graph.compile()
\end{lstlisting}

\vspace{10pt}
\textbf{Use cases:} Breaking complex agents into reusable modules, orchestrating teams of sub-agents.
\end{frame}

% ============================================================================
% SLIDE 21: Dynamic Graph Modification at Runtime
% ============================================================================
\begin{frame}[fragile]
\frametitle{Dynamic Graph Modification at Runtime}
\begin{itemize}
    \setlength{\itemsep}{1pt}
  \item \textbf{Note:} LangGraph graphs are typically compiled once
  \item For true dynamic routing, use conditional\_edges with complex logic
  \item \texttt{dspy}, \texttt{guidance}, or similar tools can modify prompts dynamically
  \item Avoid rebuilding the entire graph on each invocation
\end{itemize}

\vspace{10pt}
\begin{tipbox}[Design for Flexibility]
Use rich state and conditional edges to handle dynamic logic, rather than rebuilding graphs.
\end{tipbox}

\vspace{10pt}
\begin{lstlisting}
# Conditional edges enable dynamic routing
def dynamic_router(state: AgentState) -> str:
    # Complex logic to decide next node
    if state.get("needs_tool_call"):
        return "tool_call_node"
    return "llm_node"

graph.add_conditional_edges("think", dynamic_router, {...})
\end{lstlisting}
\end{frame}

% ============================================================================
% SLIDE 22: Map-Reduce Patterns in LangGraph
% ============================================================================
\begin{frame}[fragile]
\frametitle{Map-Reduce Patterns in LangGraph}
\begin{itemize}
    \setlength{\itemsep}{1pt}
  \item \textbf{Map:} Process multiple items in parallel (e.g., analyze documents)
  \item \textbf{Reduce:} Aggregate results (e.g., summarize)
  \item LangGraph supports concurrent node execution via state + conditional edges
\end{itemize}

\vspace{10pt}
\begin{lstlisting}
# State tracks batches
class BatchState(TypedDict):
    items: list
    results: list

def process_batch(state: BatchState) -> dict:
    """Process items (could use ThreadPoolExecutor)."""
    results = [process(item) for item in state["items"]]
    return {"results": results}

def aggregate(state: BatchState) -> dict:
    """Reduce: combine results."""
    final = sum(r["value"] for r in state["results"])
    return {"final_result": final}
\end{lstlisting}

\textit{Useful for parallel processing, batch analytics, multi-document workflows.}
\end{frame}

% ============================================================================
% SLIDE 23: Error Handling and Retry in Graph Nodes
% ============================================================================
\begin{frame}[fragile]
\frametitle{Error Handling and Retry in Graph Nodes}
\begin{lstlisting}
import tenacity

@tenacity.retry(
    stop=tenacity.stop_after_attempt(3),
    wait=tenacity.wait_exponential(multiplier=1, min=2, max=10)
)
def flaky_node(state: AgentState) -> dict:
    """Node with automatic retry."""
    # This will retry up to 3 times on failure
    result = call_unstable_service(state["input"])
    return {"output": result}

# For more control, catch errors in conditional edge
def error_handler(state: AgentState) -> str:
    if "error" in state:
        return "retry"
    return "success"

graph.add_conditional_edges("call_api", error_handler, {...})
\end{lstlisting}

\vspace{10pt}
\textit{Combine retries with conditional routing for robust workflows.}
\end{frame}

% ============================================================================
% SLIDE 24: Debugging LangGraph - Visualization
% ============================================================================
\begin{frame}[fragile]
\frametitle{Debugging LangGraph: Visualizing the Graph}
\begin{lstlisting}
# Draw graph as ASCII or Mermaid
print(compiled.get_graph().draw_ascii())

# Or get as Mermaid diagram for visual tools
mermaid_str = compiled.get_graph().to_mermaid()
print(mermaid_str)

# Inspect nodes and edges
graph_obj = compiled.get_graph()
print("Nodes:", graph_obj.nodes)
print("Edges:", graph_obj.edges)
\end{lstlisting}

\vspace{10pt}
\textbf{Debugging techniques:}
\begin{itemize}
    \setlength{\itemsep}{1pt}
  \item Use \texttt{stream()} to watch state evolution
  \item Add print statements in nodes
  \item Use \texttt{get\_state()} to inspect at interrupts
  \item Check \texttt{get\_graph().draw\_ascii()} for structure validation
\end{itemize}
\end{frame}

% ============================================================================
% SLIDE 25: Debugging - Tracing and Introspection
% ============================================================================
\begin{frame}[fragile]
\frametitle{Debugging: Tracing and Introspection}
\begin{lstlisting}
# Trace execution with debug config
config = {
    "configurable": {"thread_id": "debug_1"},
    "debug": True
}

# Stream with debug info
for step in compiled.stream(initial_state, config):
    print(f"Step: {step}")

# Get full state history (with checkpointer)
# Retrieve all checkpoints for a thread
states = compiled.get_state_history(config)
for state_snapshot in states:
    print(f"Checkpoint: {state_snapshot.checkpoint}")
    print(f"Values: {state_snapshot.values}")
\end{lstlisting}

\vspace{10pt}
\textit{Checkpointing gives you full audit trail of state evolution.}
\end{frame}

% ============================================================================
% SLIDE 26: Common Pattern - ReAct Agent as a Graph
% ============================================================================
\begin{frame}[fragile]
\frametitle{Common Pattern: ReAct Agent as a Graph}
\textbf{ReAct:} Reasoning + Acting in a loop

\vspace{10pt}
\begin{lstlisting}
class ReActState(TypedDict):
    question: str
    thoughts: str
    action: str
    observation: str
    answer: str

def think_node(state: ReActState) -> dict:
    thoughts = llm.invoke(f"Think about: {state['question']}")
    return {"thoughts": thoughts}

def act_node(state: ReActState) -> dict:
    action = llm.invoke(f"Given: {state['thoughts']}, what action?")
    obs = tools.invoke(action)
    return {"action": action, "observation": obs}

def should_continue(state: ReActState) -> str:
    if "ANSWER:" in state["observation"]:
        return "end"
    return "think"
\end{lstlisting}

Graph: START -> think -> act -> (loop back to think or END based on observation)
\end{frame}

% ============================================================================
% SLIDE 27: Common Pattern - Multi-Agent Collaboration
% ============================================================================
\begin{frame}[fragile]
\frametitle{Common Pattern: Multi-Agent Collaboration}
\begin{lstlisting}
class CollabState(TypedDict):
    topic: str
    agent_a_response: str
    agent_b_response: str
    final_answer: str

def agent_a(state: CollabState) -> dict:
    resp = agent_a_llm.invoke(state["topic"])
    return {"agent_a_response": resp}

def agent_b(state: CollabState) -> dict:
    resp = agent_b_llm.invoke(state["topic"])
    return {"agent_b_response": resp}

def synthesize(state: CollabState) -> dict:
    combined = combine(state["agent_a_response"],
                       state["agent_b_response"])
    return {"final_answer": combined}

# Graph: START -> agent_a & agent_b (parallel) -> synthesize -> END
\end{lstlisting}

\textit{Use subgraphs or parallel execution for multi-agent workflows.}
\end{frame}

% ============================================================================
% SLIDE 28: Performance Considerations
% ============================================================================
\begin{frame}
\frametitle{Performance Considerations}
\begin{itemize}
    \setlength{\itemsep}{1pt}
  \item \textbf{State size:} Keep state lean; avoid storing large objects (use references)
  \item \textbf{Checkpoint overhead:} Checkpointing after each node has cost; use in-memory for dev, persistent for prod
  \item \textbf{Streaming vs invoke:} Streaming is better for long runs; invoke waits for completion
  \item \textbf{Parallel nodes:} LangGraph doesn't parallelize by default; use Python \texttt{concurrent.futures} in nodes
  \item \textbf{Async:} Use async versions (\texttt{ainvoke}, \texttt{astream\_events}) for I/O-heavy nodes
  \item \textbf{Recursion depth:} Monitor cycles to prevent runaway loops; set max iterations if needed
\end{itemize}

\vspace{15pt}
\begin{tipbox}[Optimization]
Profile your graph with \texttt{stream()} and measure checkpoint latency in production.
\end{tipbox}
\end{frame}

% ============================================================================
% SLIDE 29: Common Pitfall 1 - State Mutation
% ============================================================================
\begin{frame}[fragile]
\frametitle{Common Pitfall 1: State Mutation}
\textbf{Pitfall:} Modifying state dict in-place leads to unexpected behavior.

\vspace{10pt}
\begin{lstlisting}
# WRONG - mutates state in place
def bad_node(state: AgentState) -> dict:
    state["messages"].append("new message")  # Mutates!
    return {}  # Returns nothing, but state changed!

# CORRECT - returns new dict
def good_node(state: AgentState) -> dict:
    new_messages = state["messages"] + ["new message"]
    return {"messages": new_messages}
\end{lstlisting}

\vspace{10pt}
\textbf{Rule:} Always return a dict with fields to update; don't mutate the input state.
\end{frame}

% ============================================================================
% SLIDE 30: Common Pitfall 2 - Missing Edges
% ============================================================================
\begin{frame}[fragile]
\frametitle{Common Pitfall 2: Missing Edges or Unreachable Nodes}
\textbf{Pitfall:} A node has no outgoing edge, causing execution to hang.

\vspace{10pt}
\begin{lstlisting}
# WRONG - "process" node has no outgoing edge
graph.add_node("process", process_fn)
graph.add_edge(START, "process")
# Missing: graph.add_edge("process", END)

# CORRECT - ensure all nodes have edges
graph.add_edge("process", END)
\end{lstlisting}

\vspace{10pt}
\textbf{Checklist:}
\begin{itemize}
    \setlength{\itemsep}{1pt}
  \item START has at least one outgoing edge
  \item Every non-END node has at least one outgoing edge
  \item Every node is reachable from START
  \item Conditional edges cover all routing function outputs
\end{itemize}
\end{frame}

% ============================================================================
% SLIDE 31: Common Pitfall 3 - Infinite Loops
% ============================================================================
\begin{frame}[fragile]
\frametitle{Common Pitfall 3: Infinite Loops}
\textbf{Pitfall:} Routing logic never reaches END.

\vspace{10pt}
\begin{lstlisting}
# WRONG - always routes back to self
def bad_route(state: AgentState) -> str:
    return "loop_node"  # Never exits!

# CORRECT - has exit condition
def good_route(state: AgentState) -> str:
    iterations = state.get("iteration", 0)
    if iterations >= 3:
        return "end"
    return "loop_node"

graph.add_conditional_edges("loop_node", good_route, {
    "loop_node": "loop_node",
    "end": END
})
\end{lstlisting}

\vspace{10pt}
\textbf{Prevention:} Add counters, check for convergence, set max recursion limits.
\end{frame}

% ============================================================================
% SLIDE 32: Common Pitfall 4 - State Type Mismatches
% ============================================================================
\begin{frame}[fragile]
\frametitle{Common Pitfall 4: State Type Mismatches}
\textbf{Pitfall:} Nodes return types not matching TypedDict schema.

\vspace{10pt}
\begin{lstlisting}
class MyState(TypedDict):
    count: int
    items: list

# WRONG - returns wrong type
def bad_node(state: MyState) -> dict:
    return {"count": "five"}  # Should be int!

# CORRECT - return correct types
def good_node(state: MyState) -> dict:
    return {"count": state["count"] + 1}
\end{lstlisting}

\vspace{10pt}
\textbf{Tip:} Use strict type checking or mypy to catch these errors early.
\end{frame}

% ============================================================================
% SLIDE 33: Cheat Sheet - Creating a Graph
% ============================================================================
\begin{frame}[fragile]
\frametitle{Cheat Sheet: Creating a Graph (1/2)}
\begin{lstlisting}
from typing_extensions import TypedDict
from langgraph.graph import StateGraph, START, END
from langgraph.checkpoint.memory import MemorySaver

# 1. Define state
class AgentState(TypedDict):
    user_input: str
    output: str

# 2. Define nodes
def process_node(state: AgentState) -> dict:
    return {"output": state["user_input"].upper()}

# 3. Create graph
graph = StateGraph(AgentState)

# 4. Add nodes
graph.add_node("process", process_node)

# 5. Add edges
graph.add_edge(START, "process")
graph.add_edge("process", END)
\end{lstlisting}
\end{frame}

% ============================================================================
% SLIDE 34: Cheat Sheet - Compiling and Invoking
% ============================================================================
\begin{frame}[fragile]
\frametitle{Cheat Sheet: Compiling and Invoking (2/2)}
\begin{lstlisting}
# 6. Compile
compiled = graph.compile(checkpointer=MemorySaver())

# 7. Invoke
result = compiled.invoke(
    {"user_input": "hello"},
    config={"configurable": {"thread_id": "1"}}
)
print(result["output"])

# 8. Stream (see intermediate steps)
for snapshot in compiled.stream(
    {"user_input": "hello"},
    config={"configurable": {"thread_id": "2"}}
):
    print(snapshot)

# 9. Get state at any point
state = compiled.get_state(config)
print(state.values)
\end{lstlisting}
\end{frame}

% ============================================================================
% SLIDE 35: Comparison - LangChain Chain vs LangGraph
% ============================================================================
\begin{frame}[fragile]
\frametitle{Comparison: LangChain Chain vs LangGraph}
\textbf{LangChain Chain (simple linear):}
\begin{lstlisting}
from langchain_core.runnables import RunnableSequence
chain = step1 | step2 | step3
result = chain.invoke({"input": "..."})
\end{lstlisting}

\vspace{10pt}
\textbf{LangGraph (stateful, looping):}
\begin{lstlisting}
graph = StateGraph(MyState)
graph.add_node("step1", step1)
graph.add_node("step2", step2)
graph.add_node("step3", step3)
graph.add_edge(START, "step1")
graph.add_conditional_edges("step2", route, {...})
compiled = graph.compile()
result = compiled.invoke({"input": "..."})
\end{lstlisting}

\vspace{10pt}
\textbf{When to use:}
\begin{itemize}
    \setlength{\itemsep}{1pt}
  \item \textbf{Chain:} Simple linear pipelines
  \item \textbf{LangGraph:} Complex, looping, interactive agents
\end{itemize}
\end{frame}

% ============================================================================
% SLIDE 36: Integration with LangChain Tools
% ============================================================================
\begin{frame}[fragile]
\frametitle{Integration with LangChain Tools}
\begin{lstlisting}
from langchain_core.tools import tool
from langgraph.prebuilt.tool_executor import ToolExecutor

@tool
def search(query: str) -> str:
    """Search the web."""
    return f"Results for {query}"

tools = [search]
tool_executor = ToolExecutor(tools)

def tool_node(state: AgentState) -> dict:
    tool_calls = state["tool_calls"]
    results = [tool_executor.invoke(tc) for tc in tool_calls]
    return {"tool_results": results}

graph.add_node("tools", tool_node)
graph.add_edge("agent", "tools")
graph.add_edge("tools", "agent")
\end{lstlisting}

\textit{LangGraph + LangChain tools = powerful agent loops.}
\end{frame}

% ============================================================================
% SLIDE 37: LLM Integration - Structured Output
% ============================================================================
\begin{frame}[fragile]
\frametitle{LLM Integration: Structured Output}
\begin{lstlisting}
from pydantic import BaseModel
from langchain_core.language_models import BaseLanguageModel

class AgentAction(BaseModel):
    action: str
    input: str

def llm_node(state: AgentState) -> dict:
    # Use structured output from LLM
    llm = ChatOpenAI(model="gpt-4").with_structured_output(
        AgentAction
    )
    action = llm.invoke([...])
    return {
        "action": action.action,
        "input": action.input
    }
\end{lstlisting}

\vspace{10pt}
\textit{Ensure LLM outputs conform to your state schema.}
\end{frame}

% ============================================================================
% SLIDE 38: Testing LangGraph Agents
% ============================================================================
\begin{frame}[fragile]
\frametitle{Testing LangGraph Agents}
\begin{lstlisting}
import pytest

def test_simple_agent():
    graph = create_agent_graph()
    compiled = graph.compile()

    result = compiled.invoke({
        "user_input": "Hello"
    })

    assert "output" in result
    assert result["output"] is not None

def test_conditional_routing():
    # Test that routing works
    result = compiled.invoke({
        "user_input": "URGENT",
        "priority": "high"
    })
    assert result["routed_to"] == "urgent_handler"
\end{lstlisting}

\vspace{10pt}
\textit{Unit test nodes independently, then test full graph flows.}
\end{frame}

% ============================================================================
% SLIDE 39: Production Deployment Considerations
% ============================================================================
\begin{frame}
\frametitle{Production Deployment Considerations}
\begin{itemize}
    \setlength{\itemsep}{1pt}
  \item \textbf{Checkpointing:} Use PostgreSQL or SQLite in production, not MemorySaver
  \item \textbf{Monitoring:} Log node execution times, errors, and state transitions
  \item \textbf{Versioning:} Version your graph definitions for reproducibility
  \item \textbf{Resource limits:} Set timeouts on LLM calls, tool execution
  \item \textbf{Async:} Use async nodes for I/O-bound operations
  \item \textbf{Error handling:} Implement retry logic and graceful degradation
  \item \textbf{State cleanup:} Archive old thread states to avoid DB bloat
  \item \textbf{API integration:} Expose graph via FastAPI, handle concurrent requests
\end{itemize}

\vspace{15pt}
\textit{LangGraph is production-ready; plan for scale and resilience.}
\end{frame}

% ============================================================================
% SLIDE 40: Self-Check Questions
% ============================================================================
\begin{frame}
\frametitle{Self-Check Questions (1/2)}
\begin{enumerate}
    \setlength{\itemsep}{1pt}
  \item What is the key difference between LangChain chains and LangGraph?
  \item How is state represented in LangGraph? What's the modern pattern?
  \item Explain the relationship between nodes, edges, and state.
  \item What does a node function return, and how is that merged into state?
  \item How do you create a loop in a LangGraph?
  \item What is the purpose of checkpointing?
  \item Describe a use case for \texttt{interrupt\_before} vs \texttt{interrupt\_after}.
\end{enumerate}
\end{frame}

% ============================================================================
% SLIDE 41: Self-Check Questions (continued)
% ============================================================================
\begin{frame}
\frametitle{Self-Check Questions (2/2)}
\begin{enumerate}
    \setlength{\itemsep}{1pt}
  \setcounter{enumi}{7}
  \item What's the difference between \texttt{invoke()}, \texttt{stream()}, and \texttt{astream\_events()}?
  \item How would you implement the ReAct (Reasoning + Acting) pattern as a LangGraph?
  \item What are three common pitfalls in LangGraph design, and how do you avoid them?
\end{enumerate}

\vspace{15pt}
\textbf{Answers should cover: graph structure, state management, cycles, checkpointing, streaming, and patterns.}
\end{frame}

% ============================================================================
% SLIDE 42: Summary
% ============================================================================
\begin{frame}
\frametitle{Summary: LangGraph Fundamentals}
\begin{itemize}
    \setlength{\itemsep}{1pt}
  \item \textbf{LangGraph:} Graph-based orchestration for stateful, looping agents
  \item \textbf{Core components:} StateGraph, nodes, edges, START/END
  \item \textbf{State:} Shared dict passed and merged through nodes
  \item \textbf{Cycles:} First-class support for loops via conditional routing
  \item \textbf{Checkpointing:} Save state at each step for interrupts, resumption, persistence
  \item \textbf{Human-in-loop:} Pause execution, modify state, resume with interrupts
  \item \textbf{Streaming:} Watch state evolution with \texttt{stream()} for long runs
  \item \textbf{Patterns:} ReAct, multi-agent, map-reduce, tool-use loops
  \item \textbf{Production-ready:} Use SQLite/PostgreSQL checkpointing, async, error handling
\end{itemize}

\vspace{15pt}
\textbf{Next:} Apply LangGraph to build complex agents with reasoning, tool use, and human feedback!
\end{frame}

% ============================================================================
% SLIDE 43: Key Takeaways
% ============================================================================
\begin{frame}
\frametitle{Key Takeaways}
\begin{center}
\Large
\textbf{Graphs $\neq$ Chains}

\vspace{15pt}
\normalsize
State + Nodes + Edges + Conditional Routing = \\
Powerful, looping, interactive agents

\vspace{20pt}
\textbf{Three Rules}
\begin{enumerate}
    \setlength{\itemsep}{1pt}
  \item Define state as TypedDict
  \item Nodes return dicts (merge, don't replace)
  \item Conditional routing enables loops
\end{enumerate}

\vspace{20pt}
\textit{Master LangGraph and unlock agent complexity.}
\end{center}
\end{frame}

% ============================================================================
% SLIDE 44: Additional Resources
% ============================================================================
\begin{frame}
\frametitle{Additional Resources}
\begin{itemize}
    \setlength{\itemsep}{1pt}
  \item \textbf{LangGraph Docs:} \url{https://langgraph.dev}
  \item \textbf{StateGraph API:} TypedDict-based state definitions
  \item \textbf{Checkpointing:} MemorySaver, SqliteSaver, PostgresSaver
  \item \textbf{Prebuilt agents:} \texttt{create\_react\_agent()} in langgraph.prebuilt
  \item \textbf{Examples:} LangChain docs, LangGraph GitHub repo
  \item \textbf{Community:} LangChain Discord for patterns and troubleshooting
\end{itemize}

\vspace{20pt}
\docref{langgraph.graph}, \docref{langgraph.checkpoint}
\end{frame}

% ============================================================================
% SLIDE 45: Contact & Questions
% ============================================================================
\begin{frame}
\frametitle{Thank You!}
\begin{center}
\Large
\textbf{Questions?}

\vspace{30pt}
\normalsize
Section 11: LangGraph Fundamentals \\
Graph-based orchestration for intelligent agents

\vspace{20pt}
\textit{Build, test, deploy, and iterate with LangGraph.}
\end{center}
\end{frame}

\end{document}
